\section{Strictness at a bounded opening price}
\label{str:section}

Bounded-price strictness was quantified in the factor-revealing
analysis of \cite[Appendix~B, Corollary~2]{CGLS2026}, following the
JMS dual-fitting framework \cite{JMS2002,JainEtAl2003}.  That work
proves the gap $1/[4(7+3T)]$ for its JMS program.  Moreover, its
fixed-size dual construction in Appendix~B.1--B.3 applies to the
relaxed scalar inequalities used below: it does not need either
ordered phase variables or the first-connection constraint.  It gives
a stronger gap than the constant we will use.  The contribution of
this section is a different elementary proof, through a suffix maximum
and two slack bounds, followed by an explicit transfer to the fixed
valid-sequence replay.  We retain a small explicit gap because that
is sufficient for the payment argument.

Write $[x]_+=\max\{x,0\}$.  The following argument uses the history
of one fixed final valid sequence from Lemma~\ref{imp:surplus-history}.
In particular, all prices, free-copy lengths, budgets, and client-removal times
below belong to that one replay.  Put $t=1-\delta$ and $K=1+v^2$.
The opened sets in the replay increase, although the algorithm that finds
the final sequence need not have this monotonicity.

\begin{lemma}[Bounded-price strictness]\label{str:strictness}
For $T>0$ define
\begin{equation}\label{str:constants}
 B=T+2,\qquad \rho=\frac{1}{200B},\qquad
 c(T)=\frac{\rho^2}{100}>0.
\end{equation}
Assume $0\le\delta<1$, $0<v\le1$, and
\begin{equation}\label{str:v-condition}
 v^2\le\frac{c(T)}{6(T+2)}.
\end{equation}
For any regular facility $i\in F$ and nonempty client set $C\subseteq D$
of positive cost $P_C=\sum_{p\in C}d(i,p)$, the history in
Lemma~\ref{imp:surplus-history} satisfies
\begin{equation}\label{str:conclusion}
 \lambda\le TP_C
 \quad\Longrightarrow\quad
 \sum_{p\in C}\alpha_p
 \le \lambda+\left(2-\frac{c(T)}2\right)P_C.
\end{equation}
The set $C$ need not be a cluster of the algorithm or of an optimal solution.
\end{lemma}

\begin{proof}
We first turn the replay inequalities into a scalar suffix inequality.
After normalizing the cluster cost, two lower bounds on its slack
rule out a small total deficit: a tail index would force a large
suffix maximum, while an early interval would then carry too much
slack.  The last step transfers the resulting gap to final budgets.

\textit{Client-removal order and a suffix inequality.}
Index the $q=|C|$ clients as $1,\ldots,q$ by their removal step in the
fixed replay, breaking ties within an atomic step arbitrarily.  This order is used
only for the present suffix proof; the final-dual proof in
Section~\ref{imp:section} sorts tied removals by final budget.  For client
$p$, let $Q_p^-$ denote the opened set immediately before its removal
step, and let $z_p$ be the phase value at that pre-step state.  Thus
clients removed in one atomic step have the same $Q_p^-$ and $z_p$.
Set
\[
 d_p=d(i,p),\qquad a_p=\frac{z_p}{K}.
\]
The imported history properties are
\begin{equation}\label{str:history-used}
 z_p\le t\,d(p,Q_p^-),\qquad
 \alpha_p\le Kz_p,
 \qquad t\,d(p,Q_p^+)\le Kz_p,
\end{equation}
where $Q_p^+$ is the immediately post-step opened set, together with
nonincreasing distances after that step.  The pre-step no-overbidding
inequality at $i$ is
\begin{equation}\label{str:no-overbid}
 \sum_{j\text{ active}}[z_p-d(i,j)]_+
 +\sum_{j\text{ inactive}}[t\,d(j,Q_p^-)-d(i,j)]_+
 \le\lambda.
\end{equation}
All sums in Equation~\eqref{str:no-overbid} initially range over $D$.

For $j<p$ define
\[
 r_{jp}=
 \begin{cases}
 t\,d(j,Q_p^-)/K,&j\text{ is already inactive before the step of }p,\\
 z_p,&j\text{ is removed in the same atomic step as }p.
 \end{cases}
\]
For fixed $j$, this quantity is nonincreasing as $p>j$ advances.
Indeed it is constant within the removal step of $j$; at the first
subsequent step, Equation~\eqref{str:history-used} gives
$t\,d(j,Q_p^-)/K\le z_j$; and thereafter the opened set only grows.
This explains why the value for a tied removal is $z_p$, even though
the scaled variable used elsewhere is $a_p$.

If $j$ is already inactive, the triangle inequality in the fixed metric
extension and Equation~\eqref{str:history-used} give
\[
 r_{jp}\ge \frac{t}{K}\bigl(d(p,Q_p^-)-d(p,j)\bigr)
 \ge a_p-\frac{t}{K}(d_p+d_j)
 \ge a_p-d_p-d_j.
\]
For a tied removal the same conclusion follows from $r_{jp}=z_p\ge a_p$.
Consequently, with
\[
 A_p=\max_{s\ge p}(a_s-d_s),
\]
monotonicity implies, for every $j<p$ and $s\ge p$,
\[
 r_{jp}\ge r_{js}\ge a_s-d_s-d_j,
 \qquad\text{and hence}\qquad r_{jp}\ge A_p-d_j.
\]
The contribution of $j<p$ in Equation~\eqref{str:no-overbid} is at least
$[r_{jp}-d_j]_+$: this is immediate for a tied removal, and follows
from $K\ge1$ for an inactive client.  Each $j\ge p$ is active in the
pre-step state of $p$, including all tied removals, and contributes at
least $[a_p-d_j]_+$.  Dropping clients outside $C$ therefore proves
\begin{equation}\label{str:suffix}
 \sum_{j\ge p}[a_p-d_j]_+
 +\sum_{j<p}[A_p-2d_j]_+\le\lambda
 \qquad(1\le p\le q).
\end{equation}

\textit{Two slack bounds and their exact sum.}
Divide all costs, phases, and budgets by $P_C$, retaining their notation.
Now $\sum_jd_j=1$ and $\lambda\le T$.  Write
$D_{p-1}=\sum_{j<p}d_j$ and define
\begin{equation}\label{str:slack-definition}
 E_p=\lambda+1+D_{p-1}+(p-1)d_p-qa_p.
\end{equation}
Expanding $[x]_+=x+[-x]_+$ in Equation~\eqref{str:suffix} yields
\begin{align}
 E_p\ &\ge (p-1)(A_p-a_p+d_p)
       +\sum_{j\ge p}[d_j-a_p]_+
       +\sum_{j<p}[2d_j-A_p]_+ \\
 &\ge (p-1)(A_p-a_p+d_p)\ge0.
 \label{str:slack-envelope}
\end{align}
Alternatively, first replace $A_p$ in Equation~\eqref{str:suffix} by
its lower bound $a_p-d_p$, and then expand.  This gives
\begin{equation}\label{str:slack-truncation}
 E_p\ge\sum_{j\ge p}[d_j-a_p]_+
           +\sum_{j<p}[2d_j+d_p-a_p]_+
      \ge\sum_j[d_j-a_p]_+.
\end{equation}
In the sum of $D_{p-1}+(p-1)d_p$ over $p$, the coefficient of $d_j$
is $(q-j)+(j-1)=q-1$.  Thus the following identity is exact:
\begin{equation}\label{str:slack-identity}
 \overline E:=\frac1q\sum_{p=1}^qE_p
  =\lambda+2-\frac1q-\sum_{p=1}^qa_p.
\end{equation}

\textit{A small deficit would concentrate neither distance nor slack.}
For a contradiction suppose
\[
 \Delta:=\lambda+2-\sum_pa_p<c,
 \qquad c=c(T).
\]
Equations~\eqref{str:slack-envelope} and~\eqref{str:slack-identity} imply
\begin{equation}\label{str:small-deficit}
 0\le\overline E=\Delta-1/q<c,
 \qquad 1/q<c,\qquad q>1/c>4/\rho.
\end{equation}
The function $F(x)=\sum_j[d_j-x]_+$ is convex and nonincreasing.
Since Equation~\eqref{str:slack-identity} also gives
$\sum_pa_p\le\lambda+2-1/q<B$, Jensen's inequality and
Equation~\eqref{str:slack-truncation} imply
\begin{equation}\label{str:jensen}
 F(B/q)\le F\left(\frac1q\sum_pa_p\right)
 \le\frac1q\sum_p F(a_p)\le\overline E<c.
\end{equation}
For every index set $J$ with $|J|\le2\rho q$, it follows that
\begin{equation}\label{str:small-set-mass}
 \sum_{j\in J}d_j
 \le |J|B/q+F(B/q)<2\rho B+c
 =\frac1{100}+c<\frac1{50}.
\end{equation}

\textit{Finding a tail index with both required bounds.}
Equation~\eqref{str:jensen} shows that fewer than $cq/B$ indices have
$d_j>2B/q$, because each such index contributes more than $B/q$ to
$F(B/q)$.  Equation~\eqref{str:small-deficit} shows that fewer than
$10cq$ indices have $E_j>1/10$.  Put $r=\lfloor\rho q\rfloor$.
The tail of exactly $r$ indices is larger than the union of these two
exceptional sets: indeed
\[
 \frac rq\ge\rho-\frac1q>\rho-c>(10+1/B)c.
\]
For the last strict inequality, $B>2$ and $\rho<1/400$ give
$(11+1/B)c<12\rho^2/100<\rho$.
Hence some $s\in\{q-r+1,\ldots,q\}$ satisfies
\begin{equation}\label{str:tail-index}
 d_s\le2B/q,\qquad E_s\le1/10.
\end{equation}
The suffix starting at $s$ has at most $r\le\rho q$ indices, so
Equation~\eqref{str:small-set-mass} gives $D_{s-1}>49/50$.
Also $(q-s+1)d_s\le r(2B/q)\le2\rho B=1/100$.
Rearranging Equation~\eqref{str:slack-definition} therefore yields
\begin{equation}\label{str:tail-large}
 q(a_s-d_s)=\lambda+1+D_{s-1}-(q-s+1)d_s-E_s
 >\lambda+1.87>\lambda+1.8.
\end{equation}

\textit{An early interval forces too much slack.}
Consider the integer interval
\[
 I=\{\lceil\rho q\rceil,\ldots,\lfloor2\rho q\rfloor\}.
\]
For $p\in I$, the prefix of length $p-1$ has at most $2\rho q$ indices.
Using $E_p,d_p\ge0$ and Equation~\eqref{str:small-set-mass} gives
\begin{equation}\label{str:early-small}
 q(a_p-d_p)=\lambda+1+D_{p-1}-(q-p+1)d_p-E_p
 <\lambda+1.02.
\end{equation}
Every $p\in I$ precedes $s$: indeed
$p\le2\rho q<(1-\rho)q+1\le q-r+1\le s$.
Thus $A_p\ge a_s-d_s$, and Equations~\eqref{str:slack-envelope},
\eqref{str:tail-large}, and~\eqref{str:early-small} imply
\[
 E_p\ge(p-1)(A_p-a_p+d_p)>\frac{0.7(p-1)}q.
\]
To check the integer rounding, set $x=\rho q>4$.
Then
\[
 |I|=\lfloor2x\rfloor-\lceil x\rceil+1>x-1>x/2,
 \qquad p-1\ge\lceil x\rceil-1\ge x-1>x/2
 \quad(p\in I).
\]
Consequently
\[
 \overline E\ge\frac1q\sum_{p\in I}E_p
 >\frac{0.7}{q^2}\frac{x}{2}\frac{x}{2}
 =0.175\rho^2>\frac{\rho^2}{100}=c,
\]
contradicting Equation~\eqref{str:small-deficit}.  We have proved
$\sum_pa_p\le\lambda+2-c$.

\textit{Returning to the final budgets.}
By Equation~\eqref{str:history-used}, $\alpha_p\le K^2a_p$.
As $v\le1$, $K^2-1=2v^2+v^4\le3v^2$, whence
\[
 \sum_p\alpha_p
 \le K^2\sum_pa_p
 \le\lambda+2-c+(K^2-1)B
 \le\lambda+2-c/2
\]
by Equation~\eqref{str:v-condition}.  Multiplication by $P_C$ proves
Equation~\eqref{str:conclusion}.
\end{proof}
